
\chapter{Duality}\label{V}

This exposé contains several theorems and compatibility results, all related
to Poincaré duality. In Section \ref{V:1}, we give the local biduality theorem
in dimension $1$ \cite[I.5.1]{sga5}, together with some computations of
duals. In Section \ref{V:2}, we give a very economical proof of Poincaré
duality for curves, taught to me by M.\ Artin. In Section \ref{V:3}, we prove
a compatibility linking two definitions of the pairing giving rise to
Poincaré duality for curves: one by cup product, the other by autoduality of
the Jacobian. Finally, in Section \ref{V:4}, we prove the compatibility
announced in the title.

Throughout the exposé, all schemes are noetherian and separated, and $n$ is
an integer invertible on every scheme under consideration.


\section{Local biduality in dimension \texorpdfstring{$1$}{1}}\label{V:1}


\subsection{}\label{V:1-1}

Let $S$ be a regular scheme, purely of dimension $1$. We shall show that the
complex with value $\dZ/n$ concentrated in degree $0$ is dualizing; that is,
for $\sK\in \ob\eD_c^b(S,\dZ/n)$ (where $(-)_c$ means constructible), if we
put $D\sK=\rHom(\sK,\dZ/n)$, then $D\sK$ is again
constructible with bounded cohomology, and the canonical morphism $\alpha$
from $\sK$ to $D D\sK$ is an isomorphism.

For $\sK$ in $\eD^-$ and arbitrary $\sL$, we have
\[
  \hom(\sK\lotimes \sL,\dZ/n) = \hom(\sL,\Hom(\sK,\dZ/n)) \text{.}
\]
The morphism from $\sK$ to $DD\sK$ is defined when $D\sK$ lies in $\eD^-$:
if $\beta:D\sK\lotimes\sK\to \dZ/n$ is the canonical pairing, it is the
morphism defined by the pairing
$\sK\lotimes D\sK=D\sK\lotimes \sK \xrightarrow\beta \dZ/n$.


\subsection{}\label{V:1-2}

Let $f:X\to S$ be a separated morphism of finite type. Put
$\sK_X=\eR f^!\dZ/n$. For $\sK\in \ob\eD^-(X,\dZ/n)$, put
$D\sK = \rHom(\sK,\sK_X)$. The adjunction between $\eR f_!$ and $\eR f^!$
gives
\[\xymatrix{
  \eR f_\ast \sK\lotimes \eR f_\ast D\sK \ar[r]
    & \eR f_! (\sK\lotimes D\sK) \ar[r]
    & \eR f_! \eR f^! \dZ/n \ar[r]
    & \dZ/n \text{.}
}\]
The symmetry of this description shows that, when $f$ is proper, the diagram
\begin{equation*}\tag{1.2.1}\label{V:eq:1-2-1}
\xymatrix{
  \eR f_\ast \sK \ar[r] \ar[d]
    & D D \eR f_\ast \sK \ar[d]^-\wr \\
  \eR f_\ast D D \sK \ar[r]^\sim
    & D\eR f_\ast D\sK
}
\end{equation*}
is commutative.

If $f$ is the inclusion of a closed point, we know that
$\eR f^!\dZ/n=\dZ/n(-1)[-2]$---that is, up to a twist and a shift, $\dZ/n$.
Over the spectrum of a field this complex is dualizing (Pontryagin duality
for $\dZ/n$-modules). By \eqref{V:eq:1-2-1}, we therefore have
$\sK\iso D D\sK$ for $\sK$ of the form $\eR f_!\sL$, and hence whenever the
support of $\underline\h^\bullet(\sK)$ is finite.


\begin{theorem_}\label{V:1-3}
Let $j:U\hookrightarrow S$ be a dense open subscheme of $S$, and let $\sF$
be a locally constant constructible sheaf of $\dZ/n$-modules on $U$. Then
$\eD j_\ast \sF = j_\ast D \sF$; that is,
$\Hom(j_\ast \sF,\dZ/n) = j_\ast\Hom(\sF,\dZ/n)$ and
$\Ext^i(j_\ast\sF,\dZ/n) = 0$ for $i>0$.
\end{theorem_}

On $U$, $\Ext^i(j_\ast\sF,\dZ/n)=0$ for $i>0$, since $\sF$ is locally
constant (and $\dZ/n$ is an injective $\dZ/n$-module), while for $i=0$ it is
the dual $\sF^\vee$ of $\sF$. One checks that
$\hom(j_\ast\sF,\dZ/n)=j_\ast\sF^\vee$, and it remains to verify the
vanishing of the $\Ext^i$ for $i>0$ at the points of $S\setminus U$. The
problem is local at these points. We are thus reduced to assuming that $S$
is a strictly local trait and that $U$ consists of its generic point $\eta$.
Let $I=\gal(\bar\eta/\eta)$. The sheaf $\sF$ is identified with the Galois
module $\sF_{\bar\eta}$, and the special fiber of $j_\ast \sF$ with
$\sF_{\bar\eta}^I$.

Let $i$ be the inclusion of the closed point $s$, and apply $D$ to the exact
sequences
\[\xymatrix{
  0 \ar[r]
    & j_! \sF \ar[r]
    & j_\ast \sF \ar[r]
    & i_\ast \sF_{\bar\eta}^I \ar[r]
    & 0 \\
  0 \ar[r]
    & j_! \sF_{\bar\eta}^I \ar[r] \ar[u]
    & \sF_{\bar\eta}^I \ar[r] \ar[u]
    & i_\ast \sF_{\bar\eta}^I \ar[r] \ar[u]
    & 0 \text{.}
}\]

We obtain a morphism of triangles
\[\xymatrix{
  i_\ast\left(\sF_{\bar\eta}^I\right)^\vee (-1)[-2] \ar[r] \ar@{=}[d]
    & D j_\ast\sF \ar[r] \ar[d]
    & \eR j_\ast \sF^\vee \ar[d] \\
  i_\ast\left(\sF_{\bar\eta}^I\right)^\vee(-1)[-2] \ar[r]
    & \left(\sF_{\bar\eta}^I\right)^\vee \ar[r]
    & \eR j_\ast \sF_{\bar\eta}^I \text{.}
}\]
The long exact sequence arising from the first row gives the vanishing of the
$\Ext^i$ for $i>2$, and, on taking the fiber at $s$, we find
\[\xymatrix@=0.7cm{
  0 \ar[r]
    & \left(\underline\h^1(D j_\ast\sF)\right)_s \ar[r]
    & \h^1\left(I,\sF_{\bar\eta}^\vee\right) \ar[r]^-\partial \ar[d]
    & \left(\sF_{\bar\eta}^I\right)^\vee(-1) \ar[r] \ar@{=}[d]
    & \left(\underline\h^2 D j_\ast\sF\right)_s \ar[r]
    & 0 \\
  & 0 \ar[r]
    & \h^1\left(I,(\sF_{\bar\eta}^I)^\vee\right) \ar[r]^-\partial
    & \left(\sF_{\bar\eta}^I\right)^\vee(-1) \ar[r]
    & 0 \text{.}
}\]
Since
$\left(\sF_{\bar\eta}^I\right)^\vee = \left(\sF_{\bar\eta}^\vee\right)_I$,
putting $M=\sF_{\bar\eta}^\vee$, it remains only to check that
\[\xymatrix{
  \h^1(I,M) \ar[r]^-\sim
    & \h^1(I,M_I) \text{.}
}\]
If $p$ is the residual characteristic exponent, then $I$ is an extension of
a group isomorphic to $\widehat\dZ_{p'} = \varprojlim_{(m,p)=1} \dZ/m$ by a
$p$-group $P$. Since the order of $P$ is prime to the order of $M$, we have
$\h^i(P,M)=0$ for $i>0$ and $(M_I)^P=(M^P)_I$. This allows us to replace $M$
by $M^P$ and $I$ by $I/P$. Finally, there is a functorial isomorphism
$\h^1(\widehat\dZ_{p'},M)\sim ($coinvariants of $\widehat\dZ_{p'}$ in $M)$,
which proves the theorem.


\begin{theorem_}\label{V:1-4}
For $\sK\in \ob\eD_c^b(S,\dZ/n)$, one has $\sK\iso D D \sK$.
\end{theorem_}

By dévissage, we reduce to the case in which $\sK$ is a constructible sheaf
$\sF$ concentrated in degree $0$, and $\sF$ either has finite support
(\ref{V:1-2}) or is of the form $j_\ast\sF_1$ as in \ref{V:1-3}. In the
second case, \ref{V:1-3} reduces us to local biduality for $\sF$ locally
constant on $U$.


\section{Poincaré duality for curves, after M.\ Artin}\label{V:2}


\subsection{}\label{V:2-1}

Let $X$ be a smooth projective curve over an algebraically closed field $k$.
Put $\sK_X=\dZ/n(1)[2]$ and, for $\sK\in\ob \eD_c^b(X,\dZ/n)$,
$D\sK=\rHom(\sK,\sK_X)$. For $M$ a $\dZ/n$-module, also put
$D M = \hom(M,\dZ/n)$, and similarly for complexes of modules. The trace
morphism $\h^0(X,\sK_X) \to \dZ/n$, or
$\eR\Gamma(X,\sK_X)\to \dZ/n$, defines a pairing
$\eR\Gamma(X,\sK)\lotimes\eR\Gamma(X,D\sK) \to \dZ/n$. Poincaré duality
between the cohomology and compactly supported cohomology of an open
subscheme $j:U\hookrightarrow X$ of the curve $X$ says that, for
$\sK=j_!\dZ/n$, this pairing identifies each factor with the dual of the
other.

I give below a proof, communicated to me by M.\ Artin, that for
$\sK\in\ob\eD_c^b(X,\dZ/n)$ this pairing is always perfect; that is, it
defines an isomorphism
\begin{equation*}\tag{2.1.1}\label{V:eq:2-1-1}
\xymatrix{
  \eR\Gamma(X,\sK) \ar[r]^-\sim
    & D \eR\Gamma(X,D\sK) \text{.}
}
\end{equation*}

For every constructible sheaf $\sF$, put
\[
  '\h^i(X,\sF) = \text{the $\dZ/n$-dual of $\h^{-i}(X,D\sF)$.}
\]
The assertion that \eqref{V:eq:2-1-1} is an isomorphism is equivalent to the
following theorem.


\begin{theorem_}\label{V:2-2}
For a constructible sheaf $\sF$ of $\dZ/n$-modules, one has
\begin{equation*}\tag{2.2.1}\label{V:eq:2-2-1}
\xymatrix{
  \h^i(X,\sF) \ar[r]^-\sim
    & '\h^i(X,\sF) \text{.}
}
\end{equation*}
\end{theorem_}


\begin{lemma_}\label{V:2-3}
Let $f:X\to Y$ be a generically étale morphism between smooth curves over
$k$. Then $\sK_X=f^\ast\sK_Y=\eR f^!\sK_Y$, with $\tr_f$ as the adjunction
morphism $\eR f_!\eR f^!\sK_Y\to \sK_Y$.
\end{lemma_}

We reduce to checking that, for $f$ finite,
$f_\ast \eR f^!\sK_Y = \rHom(f_\ast\dZ/n,\sK_Y)$ is $f_\ast\dZ/n$, with
$\tr_f$ as the adjunction morphism. The vanishing of the $\Ext^i$ for $i>0$
follows from \ref{V:1-3}, and the assertion follows.

More explicitly, $\eR f^!$ is the derived functor of the functor
$f^!:f^!\sF(U)=\hom(f_!\dZ_U,\sF)$, and \ref{V:1-3} implies the vanishing of
$\eR^i f^!\dZ/n$ for $i>0$.


\begin{corollary_}\label{V:2-4}
Let $f:X'\to X$ be a generically étale morphism of smooth projective curves
over $k$. Then $'\h^i(X,f_\ast \sF)='\h^i(X',\sF)$; this isomorphism is
functorial, and the diagram
\[\xymatrix{
  \h^i(X,f_\ast\sF) \ar[r] \ar@{=}[d]
    & '\h^i(X,f_\ast\sF) \ar@{=}[d] \\
  \h^i(X',\sF) \ar[r]
    & '\h^i(X',\sF)
}\]
is commutative.
\end{corollary_}

The isomorphism is given by \ref{V:1-2}: $D f_\ast\sF=f_\ast D\sF$.


\subsection{}\label{V:2-5}

We prove \ref{V:2-2}. If $\sF$ is $\dZ/n$ at one point, extended by $0$,
then $\h^\bullet(D\sF) = \ext^\bullet(\sF,\sK_X)=\dZ/n$ in degree $0$, and
in the pairing $\h^\bullet(\sF)\otimes\h^\bullet(D\sF) \to \dZ/n$, one has
$1\otimes 1\mapsto 1$ (\hyperref[IV]{Cycle}, \ref{IV:2-1-5}): the duality is
perfect. The case in which $\sF$ has finite support is treated in the same
way.

For arbitrary constructible $\sF$, let $\sA$ be the subsheaf of its sections
with finite support, and let $\sG=\sF/\sA$. The exact sequence
\begin{equation*}\label{V:eq:2-5-1}\tag{2.5.1}
\xymatrix{
  0 \ar[r]
    & \sA \ar[r]
    & \sF \ar[r]
    & \sG \ar[r]
    & 0
}
\end{equation*}
gives a diagram
\begin{equation*}\tag{2.5.2}\label{V:eq:2-5-2}
\xymatrix@=.5cm{
  & & 0 \ar[r]
    & \h^0(X,\sA) \ar[r] \ar[d]^-\wr
    & \h^0(X,\sF) \ar[r] \ar[d]
    & \h^0(X,\sG) \ar[r] \ar[d]
    & 0 \\
  0 \ar[r]
  & '\h^{-1}(X,\sF) \ar[r]
  & '\h^{-1}(X,\sG) \ar[r]
  & '\h^0(X,\sA) \ar[r]
  & '\h^0(X,\sF) \ar[r]
  & '\h^0(X,\sG) \ar[r]
  & 0
}
\end{equation*}
and compatible isomorphisms $\h^i(X,\sF)\iso \h^i(X,\sG)$,
$'\h^i(X,\sF) \iso '\h^i(X,\sG)$ for $i\ne 0,-1$. If $j:U\to X$ is an
open subscheme on which $\sF$ is locally constant, there is an exact sequence
\begin{equation*}\tag{2.5.3}\label{V:eq:2-5-3}
\xymatrix{
  0 \ar[r]
    & \sG \ar[r]
    & j_\ast j^\ast \sF \ar[r]
    & \sB \ar[r]
    & 0
}
\end{equation*}
with $\sB$ of finite support. Since
$'\h^i(X,j_\ast j^\ast\sF) = D \h^{2-i} (X,j_\ast(j^\ast\sF)^\vee(1))$
by \ref{V:1-3}, one has $'\h^i=0$ for $i<0$; then
\eqref{V:eq:2-5-2} and the long exact sequence defined by
\eqref{V:eq:2-5-3} show that $'\h^i$ always vanishes for $i<0$.

If $\sF=\dZ/n$, the theorem is true for $i=0$ and $2$. This expresses the
fact that, for connected $X$, the trace morphism
$\h^2(X,\dmu_n) \iso \dZ/n$ is an isomorphism. By \ref{V:2-4}, the theorem
remains true for $i=0$ and $\sF=f_\ast\dZ/n$.

For arbitrary $\sF$ again, with $\sG$ as above, there exists
$f:X'\to X$ as in \ref{V:2-4} such that $\sG$ embeds into
$f_\ast\dZ/n$:
\[\xymatrix{
  0 \ar[r]
    & \sG \ar[r]
    & f_\ast\dZ/n \ar[r]
    & \sQ \ar[r]
    & 0 \text{.}
}\]
After replacing $X'$ by an $X''/X'$, we may assume that
$\h^i(X,\sF) \to \h^i(X,f_\ast\dZ/n) = \h^i(X',\dZ/n)$ is zero for $i>0$.
Consider then
\small
\[\xymatrix@=.3cm{
  0 \ar[r]
    & \h^0(X,\sG) \ar[r] \ar[d]
    & \h^0(X',\dZ/n) \ar[r] \ar[d]^-\wr
    & \h^0(X,\sQ) \ar[r] \ar[d]
    & \h^1(X,\sG) \ar[r]^-0 \ar[d]
    & \h^1(X',\dZ/n) \ar[r] \ar[d]
    & \h^1(X,\sQ) \ar[r] \ar[d]
    & \cdots \\
  0 \ar[r]
    & '\h^0(X,\sG) \ar[r]
    & '\h^0(X',\dZ/n) \ar[r]
    & '\h^0(X,\sQ) \ar[r]
    & '\h^1(X,\sG) \ar[r]
    & '\h^1(X',\dZ/n) \ar[r]
    & '\h^1(X,\sQ) \ar[r]
    & \cdots
}\]
\normalsize

The isomorphism at $\h^0(X')$ gives
$\h^0(X,\sG)\hookrightarrow '\h^0(X,\sG)$, and the same injectivity for
$\sF$ by \eqref{V:eq:2-5-2}. Applied to $\sQ$, this gives
$\h^1(X,\sG)\hookrightarrow '\h^1(X,\sG)$. We apply this to
$\h^1(X',\dZ/n)$. This group has the same order as
$'\h^1(X',\dZ/n)=D(\h^1(X',\dZ/n)(1))$, so we obtain an isomorphism, and
$\h^1(X,\sG) \iso '\h^1(X,\sG)$. The same holds for $\sF$. For $i=2$, we
already know that $\h^2(X',\dZ/n) \iso '\h^2(X',\dZ/n)$, and proceed in the
same way, then stop for lack of combatants.


\section{Poincaré duality for curves}\label{V:3}

In this section we prove the compatibility
(\hyperref[I]{Arcata} \ref{I:eq:6-2-3-3}).


\subsection{}\label{V:3-1}

We use the notation of (\hyperref[I]{Arcata} \ref{I:6-2-3}): $\bar X$ is a
smooth, connected, projective curve over an algebraically closed field $k$,
$D$ is a reduced divisor on $\bar X$, and $0$ is a point of
$X=\bar X\setminus D$,
\[\xymatrix{
  X \ar@{^{(}->}[r]^-j
    & \bar X
    & \ar[l]_-i D \text{,}
}\]
and $\pic_D(\bar X)=\h^1(\bar X,_D\dG_m)$, where $_D\dG_m$ is the sheaf of
sections of $\dG_m$ congruent to $1\mod D$. Recall that
$\h_c^1(X,\dmu_n) = \pic_D^0(\bar X)_n$, and that $x\to \sO(x)$ defines a
canonical map $f:X\to \pic_D(\bar X)$. Put
$f_0(x)=f(x)-f(0):X\to \pic_D^0(\bar X)$.

Again write $j$ for the inclusion of $X\times X$ in $X\times \bar X$. The
diagonal $\Delta$ of $X\times X$ is closed in $X\times \bar X$; it defines a
class in
$\h_\Delta^2(X\times X,\dmu_n)=\h_\Delta^2(X\times \bar X,j_!\dmu_n)$. Let
$c$ denote its image in $\h^2(X\times \bar X,j_!\dmu_n)$.

The K\"unneth formula gives
\[
  \h^\bullet(X\times \bar X,j_!\dmu_n) = \h^\bullet(X,\dZ/n)\otimes \h^\bullet(\bar X,j_!\dmu_n) = \h^\bullet(X,\dZ/n)\otimes \h_c^\bullet(X,\dmu_n) \text{.}
\]
We shall compute the $(1,1)$-component of $c$,
\[
  c^{1,1} \in \h^1(X,\dZ/n)\otimes \h_c^1(X,\dmu_n) = \h^1(X,\h_c^1(X,\dmu_n)) = \h^1\left(X,\pic_D^0(\bar X)_n\right) \text{.}
\]

Consider the exact sequence
\[\xymatrix{
  0 \ar[r]
    & \pic_D^0(\bar X)_n \ar[r]
    & \pic_D^0(\bar X) \ar[r]^n
    & \pic_D^0(\bar X) \ar[r]
    & 0
}\]
This makes $\pic_D^0(\bar X)$ a torsor under $\pic_D^0(\bar X)_n$ over
$\pic_D^0(\bar X)$. Let $u$ be the class, in
$\h^1(X,\h_c^1(X,\dmu_n))=\h^1(X,\pic_D^0(\bar X)_n)$, of its inverse image
under $f_0$.


\begin{proposition_}\label{V:3-2}
With the preceding notation, $c^{1,1}=-u$.
\end{proposition_}

Let $e$ be the difference between the images in
$\h^2(X\times \bar X,j_!\dmu_n)$ of the cohomology classes of $\Delta$ and
$X\times\{0\}$. Since the latter is of K\"unneth type $(0,2)$, one has
$c^{1,1}=e^{1,1}$. Consider the Leray spectral sequence for
$\pr_1:X\times \bar X\to X$ and the sheaf $j_!\dmu_n$:
\[
  E_2^{pq} = \h^p(X,\eR^q {\pr_1}_\ast j_!\dmu_n) = \h^p(X,\h_c^q(X,\dmu_n)) \Rightarrow \h^{p+q}(X\times \bar X,j_!\dmu_n) \text{.}
\]
It degenerates: it is the tensor product of $\h^\bullet(X)$ with the
(trivial) Leray spectral sequence for $\bar X\to \spec(k)$ and the sheaf
$j_!\dmu_n$. Fiber by fiber, the divisor $\Delta-X\times\{0\}$ has degree
$0$, so the image of $e$ in $E^{02}$ is zero. We may therefore speak of its
image $\bar e$ in $E^{1,1}$, and we must show that $\bar e = -u$.

Let $\sG$ on $X\times\bar X$ be the subsheaf of $\dG_m$ consisting of local
sections whose restriction to the closed subscheme $X\times D$ is $1$.
There is again an exact sequence
\begin{equation*}\tag{3.2.1}\label{V:eq:3-2-1}
\xymatrix{
  0 \ar[r]
    & j_!\dmu_n \ar[r]
    & \sG \ar[r]
    & \sG \ar[r]
    & 0
}
\end{equation*}
and $e$ is the image under $\partial$ of the class
$e_1\in \h^1(X\times\bar X,\sG)$ of the invertible sheaf
$\sO(\Delta-X\times\{0\})$, trivialized by $1$ on $X\times D$.

Applying $\eR\pr_{1\ast}$ to the distinguished triangle defined by
\eqref{V:eq:3-2-1} gives a distinguished triangle
\begin{equation*}\tag{3.2.2}\label{V:eq:3-2-2}
\xymatrix{
  & \ar[r]^-\partial
    & \eR \pr_{1\ast} j_!\dmu_n \ar[r]
    & \eR \pr_{1\ast}\sG \ar[r]
    & \eR \pr_{1\ast}\sG \ar[r]^-\partial
    &
}
\end{equation*}
and the diagram
\[\xymatrix{
  \h^1(X\times \bar X,\sG) \ar[r]^-\partial \ar@{=}[d]^-{(1)}
    & \h^2(X\times \bar X,j_!\dmu_n) \ar@{=}[d] \\
  \h^1(X,\eR\pr_{1\ast} \sG) \ar[r]^-\partial
    & \h^2(X,\eR\pr_{1\ast} j_!\dmu_n)
}\]
is commutative. Our problem is thus to identify the image in
$E^{11}=\h^1(X,\h_c^1(X,\dmu_n))$ (the image in $E^{02}$ being zero) of the
image under
$\partial_{\text{\eqref{V:eq:3-2-2}}}:\h^1(X,\eR\pr_{1\ast}\sG) \to \h^2(X,\eR\pr_{1\ast} j_!\dmu_n)$
of the class, still denoted by $e_1$, corresponding to $e_1$ under (1).

The sheaf $\eR^1\pr_{1\ast}\sG$ is defined by the group scheme over $X$
obtained by pulling back the algebraic group $\pic_D(\bar X)$ over
$\spec(k)$, and the image of $e_1$ in
$\h^0(X,\eR^1\pr_{1\ast}\sG) = \hom(X,\pic_D(\bar X))$ is $f_0$.

Represent the triangle \eqref{V:eq:3-2-2} by a short exact sequence of
complexes of sheaves:
\[\xymatrix{
  0 \ar[r]
    & \sA \ar[r]
    & \sB \ar[r]
    & \sC \ar[r]
    & 0 \text{.}
}\]
Let $\sB'$ be the subsheaf of
\[
  \tau_{\leqslant 1} \sB : \cdots \sB^0 \to \ker(d) \to 0 \cdots
\]
obtained by replacing the degree-one term $\ker(d)$ by the inverse image, under
$\ker(d)\to \underline\h^1(\sB) = \pic_D(\bar X)$ over $X$, of
$\pic_D^0(\bar X)$ over $X$. Let $\sA'=\sA\cap \sB'$ and let $\sC'$ be the
image of $\sB'$. We have $\sA'=\tau_{\leqslant 1}(\sA)$, and $\sC'$ is
obtained from $\sC$ in the same way as $\sB'$ from $\sB$. Hence there is a
commutative diagram of short exact sequences
\begin{equation*}\tag{3.2.3}\label{V:eq:3-2-3}
\xymatrix{
  0 \ar[r]
    & \sA \ar[r]
    & \sB \ar[r]
    & \sC \ar[r]
    & 0 \\
  0 \ar[r]
    & \tau_{\leqslant 1} \sA \ar[r] \ar[u] \ar[d]
    & \sB' \ar[r] \ar[u] \ar[d]
    & \sC' \ar[r] \ar[u] \ar[d]
    & 0 \\
  0 \ar[r]
    & \h_c^1(X,\dmu_n)[-1] \ar[r]
    & \pic_D^0(\bar X)[-1] \ar[r]
    & \pic_D^0(\bar X)[-1] \ar[r]
    & 0
}
\end{equation*}
where the last row is to be viewed as an exact sequence of complexes of
sheaves concentrated in degree $1$ on $X$.
\begin{enumerate}[a)]
  \item $e\in \hh^2(X,\sA)$ comes from
    $\widetilde e\in \hh^2(X,\tau_{\leqslant 1}\sA)$ (because its image in
    $E^{02}$ is zero). We seek the image $\bar e$ of $\widetilde e$ in
    $\hh^2(X,\h_c^1(X,\dmu_n)[-1]) = \h^1(X,\h_c^1(X,\dmu_n))$.
  \item $e_1\in \hh^1(X,\sC)$ comes from
    $\widetilde e_1\in \hh^1(X,\sC')$, because its image $f_0$ in
    $\h^0(X,\pic_D(\bar X))$ lies in $\h^0(X,\pic_D^0(\bar X))$. We may
    take $\widetilde e=\partial \widetilde e_1$.
  \item Thus $\bar e = \partial f_0$, where $\partial$ is defined by
    \eqref{V:eq:3-2-2}, and where we use the usual isomorphisms
    $\hh^2(X,\sF[-1]) = \h^1(X,\sF)$. These isomorphisms anticommute with
    $\partial$, so that $\bar e=-\partial f_0$, for $\partial$ defined by
    the exact sequence
    \[\xymatrix{
      0 \ar[r]
        & \h_c^1(X,\dmu_n) \ar[r]
        & \pic_D^0(\bar X) \ar[r]
        & \pic_D^0(\bar X) \ar[r]
        & 0
    }\]
    The conclusion follows from (\hyperref[IV]{Cycles}, \ref{IV:1-1-1}).
\end{enumerate}


\subsection{}\label{V:3-3}

For every homomorphism $\varphi:\h_c^1(X,\dmu_n)\to \dZ/n$, let
$\varphi(u)\in \h^1(X,\dZ/n)$ be the image of $u$ under
$\varphi:\h^1(X,\h_c^1(X,\dmu_n)) \to \h^1(X,\dZ/n)$. If
$x\in \h_c^1(X,\dmu_n)$, the cup product $\varphi(u)\smallsmile x$ lies in
$\h_c^2(X,\dmu_n)$. We shall prove the following.


\begin{proposition_}\label{V:3-4}
$\tr(\varphi(u)\smallsmile x) = \varphi(x)$.
\end{proposition_}

This identity follows, after applying $\varphi$, from
\begin{equation*}\tag{3.4.1}\label{V:eq:3-4-1}
  \tr(u\smallsmile x) = x
\end{equation*}
where $\tr$ now denotes the map
\[\xymatrix{
  \tr:\h_c^2(X,\h_c^1(X,\dmu_n)\otimes\dmu_n) \ar[r]
    & \h_c^1(X,\dmu_n) \text{,}
}\]
so that $x\mapsto \tr(u\smallsmile x)$ is the composite
\small
\begin{equation*}\tag{3.4.2}\label{V:eq:3-4-2}
\begin{aligned}
&\xymatrix{
  \h_c^1(X,\dmu_n) \ar[r]^-{u\otimes x}
    & \h^1(X,\h_c^1(X,\dmu_n)) \otimes \h_c^1(X,\dmu_n)
    & \ar[l]_-\sim \h^1(X,\dZ/n)\otimes \h_c^1(X,\dmu_n) \otimes \h_c^1(X,\dmu_n)} \\ &\xymatrix{
    \ar[r]^-{(1)}
      & \h_c^2(X,\dmu_n) \otimes \h_c^1(X,\dmu_n) \ar[r]^-{\tr\otimes\operatorname{id}}
      & \h_c^1(X,\dmu_n)
}
\end{aligned}
\end{equation*}
\normalsize
where (1) is given by
$a\otimes b\otimes c\mapsto (a\smallsmile c)\otimes b$.

We express this morphism in terms of the trace morphism
$\tr:\h_c^3(X\times X,\dmu_n^{\otimes 2}) \to \h_c^2(X,\dmu_n)$ relative to
the second projection. For $x\in \h_c^3$, $\tr(x)$ is computed as follows:
among its K\"unneth components, retain only the component of type $(2,1)$,
and apply $\tr:\h_c^2(X,\dmu_n) \to \dZ/n$.

Let $u'$ be the image of $u$ under the map
\begin{align*}
&\xymatrix{
  \h^1(X,\h_c^1(X,\dmu_n))
    & \ar[l]_-\sim \h^1(X,\dZ/n) \otimes \h_c^1(X,\dmu_n) = \h^1(X,\dZ/n)\otimes \h^1(\bar X,j_!\dmu_n)
} \\
&\xymatrix{
  \ar[r]
    & \h^2(X\times \bar X,\dZ/n\boxtimes j_!\dmu_n) \text{.}
}
\end{align*}
In view of the relation between the K\"unneth formula and cup product, one
has
\[
  \tr(u\smallsmile x) = -\tr(u'\smallsmile \pr_1^\ast x) \text{.}
\]
The second cup product is the one in
(\hyperref[IV]{Cycle} \ref{IV:1-2-4}):
\[\xymatrix{
  \h^2(X\times \bar X,\dZ/n\boxtimes j_!\dmu_n) \otimes \h^1(\bar X\times \bar X,j_!\dmu_n\boxtimes \dZ/n) \ar[r]
    & \h^3(\bar X\times \bar X,j_!\dmu_n \boxtimes j_!\dmu_n) \text{.}
}\]
The minus sign comes from permuting two symbols of degree $1$ in (1). Since
$c\smallsmile\pr_1^\ast x$ and $c^{1,1}\smallsmile \pr_1^\ast x$ have the
same $(2,1)$-component in the K\"unneth decomposition, and since
$u'=-c^{1,1}$, one has
\[
  \tr(u\smallsmile x) = \tr(c\smallsmile \pr_1^\ast x) \text{.}
\]
The cup product on the right can now be interpreted as the cup product
(\hyperref[IV]{Cycle} \ref{IV:1-2-5}) of
$\cl(\Delta)\in \h_\Delta^2(X\times X,\dmu_n)$ with
$\pr_1^\ast x\in \h_c^1(\Delta,\dmu_n)$. On $\Delta$, one has
$\pr_1=\pr_2$, whence
\[
  \tr(c\smallsmile \pr_1^\ast x) = \tr(\cl(\Delta)\smallsmile \pr_1^\ast x) = \tr(\cl(\Delta)\smallsmile \pr_2^\ast x) = \tr(\cl(\Delta))\smallsmile x = x \text{.}
\]


\section{Compatibility with SGA 4 XVIII 3.1.10.3}\label{V:4}

In \cite[XVIII]{sga4}, overcome by discouragement, I did not verify the
compatibility announced in the title. Although it turns out to be no more
than smoke, I feel bound to remedy that omission here. For $f:X\to S$
separated and of finite type, the point was to verify that the adjunction
morphism $\eR f_!\eR f^!\sL\to \sL$ is compatible with every étale
localization $k:V\to S$. Using $k$ and $f$ for several arrows, as below, we
must verify commutativity:
\[\xymatrix{
  X_V \ar[r]^-k \ar[d]^-f
    & X \ar[d]^-f \\
  V \ar[r]^-k
    & S
} \qquad
\xymatrix{
  k^\ast \eR f_! \eR f^!\sL \ar[r]^-{k^\ast(\text{adj})} \ar[d]^-\sim
    & k^\ast \sL \ar@{=}[d] \\
  \eR f_! \eR f^! k^\ast \sL \ar[r]^-{\text{adj}}
    & k^\ast \sL \text{.}
}\]
At the level of complexes, the required commutativity is written (loc. cit.)
\[\xymatrix{
  k^\ast f_!^\bullet f^!_\bullet \sL \ar[r] \ar[d]^-{(1)}
    & k^\ast \sL \ar@{=}[d] \\
  f_!^\bullet f^!_\bullet k^\ast \sL \ar[r]
    & k^\ast \sL \text{.}
}\]
It is checked component by component, reducing the problem to a
compatibility for $\sL$ a sheaf and for each pair of adjoint functors
$(f_i^!,f_!^i)$, denoted simply by $f^!$ and $f_!$. The arrow (1) is defined
using the isomorphism $(2):k^\ast f_!=f_! k^\ast$, and a morphism
$(3):k^\ast f^! \to f^! k^\ast$ deduced from it. In loc. cit., one first
introduces a morphism $(4):k_! f_!\to f_! k_!$, deduced from $(2)$ by the
adjunction of $k^\ast=k^!$ and
$k_!:k_! f_!\to k_! f_! k^\ast k_!\to k_! k^\ast f_! k_! \to f_! k_!$, and
then deduces $(3)$ from $(4)$ by adjunction. Equivalently, one may deduce
$(3)$ from $(2)$ by the adjunction of $f_!$ and
$f^!:k^\ast \to f^! f_! k^\ast f^!\to f^! k^\ast f_! f^! \to f^! k^\ast$.

Writing $k$ for $k^\ast$ and omitting $\sL$, the required commutativity is
that of the outer boundary of
\[\xymatrix@=.5cm{
  k f_! f^! \ar[r] \ar@{=}[d]
    & f_! k f^! \ar[r]
    & f_! f^! f_! k f^! \ar[r]
    & f_! f^! k f_! f^! \ar@{=}[d] \\
  k f_! f^! \ar[d]
    & & & \ar[lll] f_! f^! k f_! f^! \ar[d] \\
  k
    & & & \ar[lll] f_! f^! k \text{.}
}\]
In the first row, we find a homomorphism
$k f_! f^! \to f_! f^! k f_! f^!$, defined because $k f_! f^!$ is of the
form $f_! X$ (with $X=k f^!$). It admits the adjunction morphism in the
second row as a retraction, and the commutativity of the lower square is the
functoriality of this adjunction morphism.
